\begin{thebibliography}{Stacks}

\bibitem[And92]{Andre1992}
Y.~Andr\'e,
\emph{Une remarque \`a propos des cycles de Hodge de type CM},
in S.~David (ed.), \emph{S\'eminaire de Th\'eorie des Nombres,
Paris, 1989--90}, Progress in Math. \textbf{102},
Birkh\"auser Boston, Boston, MA, 1992, 1--7.

\bibitem[And96]{Andre1996}
Y.~Andr\'e,
\emph{Pour une th\'eorie inconditionnelle des motifs},
Publ. Math. Inst. Hautes \'{E}tudes Sci. \textbf{83} (1996), 5--49.
\href{https://doi.org/10.1007/BF02698643}{doi:10.1007/BF02698643}.
\href{https://www.numdam.org/item/10.1007/BF02698643.pdf}{Original article}.

\bibitem[Bla94]{Blasius1994}
D.~Blasius,
\emph{A $p$-adic property of Hodge classes on abelian varieties},
in \emph{Motives} (Seattle, WA, 1991), Part~2,
U.~Jannsen, S.~Kleiman, and J.-P.~Serre (eds.),
Proc. Sympos. Pure Math. \textbf{55}, Part~2,
American Mathematical Society, Providence, RI, 1994, 293--308.
\href{https://doi.org/10.1090/pspum/055.2/1265557}{doi:10.1090/pspum/055.2/1265557}.
\href{https://archive.mpim-bonn.mpg.de/id/eprint/3859/1/preprint_1992_59.pdf}{Author's preprint}.

\bibitem[CF00]{ColmezFontaine2000}
P.~Colmez and J.-M.~Fontaine,
\emph{Construction des repr\'esentations $p$-adiques semi-stables},
Invent. Math. \textbf{140} (2000), no.~1, 1--43.
\href{https://doi.org/10.1007/s002220000042}{doi:10.1007/s002220000042}.
\href{https://www.imo.universite-paris-saclay.fr/~fontaine/CF.pdf}{Author's version}.

\bibitem[Con03]{Conrad2003}
B.~Conrad,
\emph{Minimal models for elliptic curves}, notes, 20 October 2003.
\href{https://math.stanford.edu/~conrad/vigregroup/vigre03/minimalmodel.pdf}{Author's notes}.

\bibitem[Del82]{Deligne1982}
P.~Deligne,
\emph{Hodge cycles on abelian varieties} (notes by J.~S.~Milne),
in P.~Deligne, J.~S.~Milne, A.~Ogus, and K.-y.~Shih,
\emph{Hodge Cycles, Motives, and Shimura Varieties},
Lecture Notes in Math. \textbf{900}, Springer-Verlag,
Berlin--New York, 1982, 9--100.
\href{https://www.jmilne.org/math/Documents/Deligne82.pdf}{Corrected online version}.

\bibitem[FR05]{FontaineRapoport2005}
J.-M.~Fontaine and M.~Rapoport,
\emph{Existence de filtrations admissibles sur des isocristaux},
Bull. Soc. Math. France \textbf{133} (2005), no.~1, 73--86.
\href{https://doi.org/10.24033/bsmf.2479}{doi:10.24033/bsmf.2479}.
\href{https://www.numdam.org/article/BSMF_2005__133_1_73_0.pdf}{Published version}.

\bibitem[GLX26]{GleasonLimXu2026}
I.~Gleason, D.~G.~Lim, and Y.~Xu,
\emph{The connected components of affine Deligne--Lusztig varieties},
Invent. Math. \textbf{243} (2026), 805--861.
\href{https://doi.org/10.1007/s00222-025-01386-1}{doi:10.1007/s00222-025-01386-1}.
\href{https://arxiv.org/abs/2208.07195v3}{arXiv:2208.07195v3}
(10 November 2025); preprint locators refer to
\href{https://arxiv.org/pdf/2208.07195v3}{this version}.

\bibitem[Gro71]{Grothendieck1970}
A.~Grothendieck,
\emph{Groupes de Barsotti-Tate et cristaux},
in \emph{Actes du Congr\`es International des Math\'ematiciens}
(Nice, 1970), Tome~1, Gauthier-Villars, Paris, 1971, 431--436.
\href{https://webusers.imj-prg.fr/~leila.schneps/grothendieckcircle/GroupesB-T.pdf}{Original article}.

\bibitem[Har08]{Hartl2008}
U.~Hartl,
\emph{On period spaces for $p$-divisible groups},
C. R. Math. Acad. Sci. Paris \textbf{346} (2008), nos.~21--22, 1123--1128.
\href{https://doi.org/10.1016/j.crma.2008.07.003}{doi:10.1016/j.crma.2008.07.003}.
\href{https://comptes-rendus.academie-sciences.fr/mathematique/item/10.1016/j.crma.2008.07.003.pdf}{Published version}.

\bibitem[Ill15]{IllusiePisa}
L.~Illusie,
\emph{Grothendieck at Pisa: crystals and Barsotti--Tate groups},
in U.~Zannier (ed.), \emph{Colloquium De Giorgi 2013 and 2014},
Publications of the Scuola Normale Superiore \textbf{5},
Edizioni della Normale, Pisa, 2015, 79--107.
\href{https://doi.org/10.1007/978-88-7642-515-8_5}{doi:10.1007/978-88-7642-515-8\_5}.
\href{https://www.imo.universite-paris-saclay.fr/~luc.illusie/Illusie-Pisa5.pdf}{Author's version}.

\bibitem[Kim12]{Kim2012}
W.~Kim,
\emph{The classification of $p$-divisible groups over $2$-adic discrete
valuation rings},
Math. Res. Lett. \textbf{19} (2012), no.~1, 121--141.
\href{https://doi.org/10.4310/MRL.2012.v19.n1.a10}{doi:10.4310/MRL.2012.v19.n1.a10}.
\href{https://arxiv.org/abs/1007.1904v3}{arXiv:1007.1904v3}
(30 December 2011); preprint locators refer to
\href{https://arxiv.org/pdf/1007.1904v3}{this version}.

\bibitem[Kis06]{Kisin2006}
M.~Kisin,
\emph{Crystalline representations and $F$-crystals},
in V.~Ginzburg (ed.), \emph{Algebraic Geometry and Number Theory:
In Honor of Vladimir Drinfeld's 50th Birthday},
Progress in Math. \textbf{253}, Birkh\"auser Boston, Boston, MA,
2006, 459--496.
\href{https://doi.org/10.1007/978-0-8176-4532-8_7}{doi:10.1007/978-0-8176-4532-8\_7}.

\bibitem[Kis17]{Kisin2017}
M.~Kisin,
\emph{Mod $p$ points on Shimura varieties of abelian type},
J. Amer. Math. Soc. \textbf{30} (2017), no.~3, 819--914.
\href{https://people.math.harvard.edu/~kisin/dvifiles/lr.pdf}{Author's version used for all locators}.

\bibitem[KZ25a]{KisinZhouEarlier}
M.~Kisin and R.~Zhou,
\emph{Independence of $\ell$ for Frobenius conjugacy classes attached to
abelian varieties},
Ann. of Math. (2) \textbf{202} (2025), no.~3, 1077--1156.
\href{https://doi.org/10.4007/annals.2025.202.3.3}{doi:10.4007/annals.2025.202.3.3}.
\href{https://arxiv.org/abs/2103.09945v2}{arXiv:2103.09945v2} (7 October 2024).
All locators refer to
\href{https://arxiv.org/pdf/2103.09945v2}{arXiv:2103.09945v2}.

\bibitem[KZ25b]{KisinZhou2025}
M.~Kisin and R.~Zhou,
\emph{Strongly compatible systems associated to semistable abelian varieties},
preprint, 2025,
\href{https://arxiv.org/abs/2505.02165v1}{arXiv:2505.02165v1}
(deposited 4 May 2025).
\href{https://arxiv.org/pdf/2505.02165v1}{PDF version used for all locators},
dated 10 June 2025.

\bibitem[Mar26]{Markman2026}
E.~Markman,
\emph{Secant sheaves and Weil classes on abelian varieties},
preprint, contribution to the ICM 2026,
\href{https://arxiv.org/abs/2509.23403v2}{arXiv:2509.23403v2}
(11 February 2026).
\href{https://arxiv.org/pdf/2509.23403v2}{Version used for all locators},
dated 13 February 2026.

\bibitem[Mil86]{MilneAV}
J.~S.~Milne,
\emph{Abelian varieties},
in G.~Cornell and J.~H.~Silverman (eds.), \emph{Arithmetic Geometry}
(Storrs, CT, 1984), Springer-Verlag, New York, 1986, 103--150.
\href{https://doi.org/10.1007/978-1-4613-8655-1_5}{doi:10.1007/978-1-4613-8655-1\_5}.
\href{https://www.jmilne.org/math/xnotes/AVs.pdf}{Corrected author's version},
2 January 2022; original theorem numbering retained.

\bibitem[Mil99]{Milne1999}
J.~S.~Milne,
\emph{Lefschetz classes on abelian varieties},
Duke Math. J. \textbf{96} (1999), no.~3, 639--675.
\href{https://www.jmilne.org/math/articles/1999aP.pdf}{Author's version};
\href{https://www.jmilne.org/math/articles/1999a.html}{author's correction to
the proof of Theorem~5.9}.

\bibitem[Mil00]{Milne2000}
J.~S.~Milne,
\emph{Towards a proof of the conjecture of Langlands and Rapoport},
notes for a talk at the Institut Henri Poincar\'e, 28 April 2000,
version dated 30 April 2000.
\href{https://www.jmilne.org/math/articles/2000.pdf}{Author's notes}.

\bibitem[Mil05]{MilneShimura}
J.~S.~Milne,
\emph{Introduction to Shimura varieties},
in \emph{Harmonic Analysis, the Trace Formula, and Shimura Varieties},
Clay Math. Proc. \textbf{4}, American Mathematical Society,
Providence, RI, 2005, 265--378.
\href{https://www.jmilne.org/math/xnotes/svi.pdf}{Author's revision},
16 September 2017; locators refer to this revision.

\bibitem[Mil08a]{MilneAIM2007}
J.~S.~Milne,
\emph{The Tate conjecture over finite fields (AIM talk)},
notes for the American Institute of Mathematics workshop, July 2007,
version~2.2, 7 May 2008.
\href{https://www.jmilne.org/math/articles/2007e.pdf}{Version used for all locators}.

\bibitem[Mil08b]{MilneAVNotes}
J.~S.~Milne,
\emph{Abelian Varieties}, course notes, version~2.00, 2008, vi+166~pp.
\href{https://www.jmilne.org/math/CourseNotes/AVc.pdf}{Author's course notes}.

\bibitem[Mil09]{Milne2009}
J.~S.~Milne,
\emph{Rational Tate classes},
Mosc. Math. J. \textbf{9} (2009), no.~1, 111--141.
\href{https://www.jmilne.org/math/articles/2009b.pdf}{Author's version};
page numbers refer to this version.

\bibitem[Mil13]{MilneLEC}
J.~S.~Milne,
\emph{Lectures on \'{E}tale Cohomology},
course notes, version~2.21, 22 March 2013.
\href{https://www.jmilne.org/math/CourseNotes/LEC.pdf}{Author's notes}.

\bibitem[MilANT]{MilneANT}
J.~S.~Milne,
\emph{Algebraic Number Theory}, course notes, version~3.08, 19 July 2020.
\href{https://www.jmilne.org/math/CourseNotes/ANTc.pdf}{Author's course notes}.
Locators refer to this version.

\bibitem[Mil20]{Milne2020}
J.~S.~Milne,
\emph{Hodge classes on abelian varieties},
preprint, 2020,
\href{https://arxiv.org/abs/2010.08857v2}{arXiv:2010.08857v2}
(11 November 2020).
\href{https://arxiv.org/pdf/2010.08857v2}{Version used for all locators}.

\bibitem[Mil22]{MilneAG2022}
J.~S.~Milne,
\emph{Algebraic Groups: The Theory of Group Schemes of Finite Type over a Field},
Cambridge Studies in Advanced Mathematics \textbf{170},
Cambridge University Press, Cambridge, 2017; corrected reprint, 2022.
\href{https://doi.org/10.1017/9781316711736}{doi:10.1017/9781316711736}.
\href{https://www.jmilne.org/math/Books/iAG2022.pdf}{Corrected 2022 version used for all locators}.

\bibitem[Mil25]{Milne2025}
J.~S.~Milne,
\emph{Abelian motives and Shimura varieties in nonzero characteristic},
preprint, version~2.2, 2 November 2025.
\href{https://www.jmilne.org/math/articles/SVp.pdf}{Author's version used for all locators}.

\bibitem[Moo98]{Moonen1998}
B.~Moonen,
\emph{Models of Shimura varieties in mixed characteristics},
in A.~J.~Scholl and R.~L.~Taylor (eds.),
\emph{Galois Representations in Arithmetic Algebraic Geometry}
(Durham, 1996), London Math. Soc. Lecture Note Ser. \textbf{254},
Cambridge University Press, Cambridge, 1998, 267--350.
\href{https://doi.org/10.1017/CBO9780511662010.008}{doi:10.1017/CBO9780511662010.008}.
\href{https://www.math.ru.nl/personal/bmoonen/Papers/SMCfinal.pdf}{Author's version}.

\bibitem[OAI26]{OpenAICMHodge2026}
OpenAI,
\emph{The rational Hodge conjecture for CM abelian varieties},\newline
OpenAI Math Release preprint
\href{https://github.com/openai/math/blob/main/preprints/The-rational-Hodge-conjecture-for-CM-abelian-varieties-September-30-2026/paper.pdf}{OAI:The-rational-Hodge-conjecture-for-CM-abelian-varieties-September-30-2026},
2026.

\bibitem[PR24]{PappasRapoport2024}
G.~Pappas and M.~Rapoport,
\emph{$p$-adic shtukas and the theory of global and local Shimura varieties},
Camb. J. Math. \textbf{12} (2024), no.~1, 1--164.
\href{https://doi.org/10.4310/CJM.2024.v12.n1.a1}{doi:10.4310/CJM.2024.v12.n1.a1}.
\href{https://arxiv.org/abs/2106.08270v6}{arXiv:2106.08270v6}
(26 October 2023); preprint page numbers refer to
\href{https://arxiv.org/pdf/2106.08270v6}{this version}.

\bibitem[Stacks]{StacksProject}
The Stacks Project Authors,
\emph{The Stacks Project},
\href{https://stacks.math.columbia.edu/tag/0A42}{Lemma~30.28.3, Tag~0A42},
accessed 23 September 2026.

\bibitem[Tat67]{Tate1967}
J.~T.~Tate,
\emph{$p$-Divisible Groups},
in T.~A.~Springer (ed.), \emph{Proceedings of a Conference on Local Fields}
(Driebergen, 1966), Springer-Verlag, Berlin--Heidelberg, 1967, 158--183.
\href{https://doi.org/10.1007/978-3-642-87942-5_12}{doi:10.1007/978-3-642-87942-5\_12}.
\href{https://www.math.purdue.edu/~tongliu/teaching/598/p-divisible.pdf}{Original article}.

\bibitem[Win97]{Wintenberger1997}
J.-P.~Wintenberger,
\emph{Propri\'et\'es du groupe tannakien des structures de Hodge $p$-adiques
et torseur entre cohomologies cristalline et \'etale},
Ann. Inst. Fourier (Grenoble) \textbf{47} (1997), no.~5, 1289--1334.
\href{https://doi.org/10.5802/aif.1600}{doi:10.5802/aif.1600}.
\href{https://www.numdam.org/item/10.5802/aif.1600.pdf}{Original article}.

\end{thebibliography}
